\section{About the Examples}
\label{sec:examples}

The examples use the Philippine CLEWS case, comparing a national energy-policy scenario
with a baseline, linked to a calibration of OG-Core for the Philippines with eight
industries. The model solutions were saved in July 2026. Each effect is measured against a
control: the same model solved as a scenario but without any channel adjustment. Solving a
scenario introduces small differences from the original baseline, mainly in consumption in
the first few years; measuring against the control keeps them out of the channel effects.

The magnitudes depend on the calibration of both models; the Philippine calibration, for
instance, sets no minimum level of electricity use for households. Three limits apply to
these examples. The economy's demand and cost of capital were computed but not yet passed back
to CLEWS for a new solution, so the examples show one exchange rather than the full
two-way loop. The investment incentive was not part of the saved runs. And the conversion
between the monetary units of the two models is still being calibrated; this affects the
size of the carbon-price and investment examples, but not the electricity-cost or health
examples, which use relative changes.

\pending{note for the author, not for the reader: the health results quoted in
sections 3 and 5 come from the corrected run's record. The remaining channel numbers
and all figures still reflect the earlier frozen outputs and await a re-freeze on the
current run record --- whose electricity-cost rise is larger and sustained (no
mid-2030s relief) and whose combined long-run output effect is about 1.0 percent,
not 0.14.}
